\appendix
\section*{Appendix R8: NDH Resonance Visualization Formalism}

\subsection*{R8.0 Overview}

This appendix formalizes the geometric and algebraic mechanisms by which
the NDH 42D interpretation manifold $X_{42}$ produces stable, coherent,
oscillatory structures when projected into ternary visualization engines.
These resonance artifacts arise from the fixed-point invariance of the
interpretation operator $I_p$ and the ternary-closed structure of the
NDH operator algebra.

\subsection*{R8.1 Resonance Fields on the 42D Manifold}

Let $X_{42}$ be the 42-dimensional interpretation manifold defined in
Appendix R4.

Let $\mathfrak{X}(X_{42})$ denote the space of smooth vector fields on
$X_{42}$.

Define the resonance-field family:
\[
\mathcal{R} = \{ r_{\alpha} \in \mathfrak{X}(X_{42}) \mid \alpha \in A \},
\]
where each $r_{\alpha}$ satisfies:

\begin{enumerate}
    \item smoothness: $r_{\alpha}$ is $C^\infty$,
    \item boundedness: $\|r_{\alpha}(x)\| < K$ for some constant $K$,
    \item ternary closure:
    \[
    r_{\alpha} \otimes r_{\beta} \otimes r_{\gamma} \in \mathcal{R}.
    \]
\end{enumerate}

These fields generate resonance trajectories on $X_{42}$.

\subsection*{R8.2 Resonance Flow}

Define the resonance flow $\Phi_{\alpha} : X_{42} \times \mathbb{R} \to X_{42}$ by:
\[
\frac{d}{dt} \Phi_{\alpha}(x,t) = r_{\alpha}(\Phi_{\alpha}(x,t)),
\qquad \Phi_{\alpha}(x,0) = x.
\]

The flow is well-defined because $r_{\alpha}$ is smooth and bounded.

\subsection*{R8.3 Projection into Ternary Visualization Space}

Let $V_3$ be a ternary visualization space (e.g., a 3D or 4D rendering
engine with ternary tensor collapse).

Define a projection:
\[
\Pi : X_{42} \to V_3
\]
satisfying:

\begin{enumerate}
    \item smoothness,
    \item surjectivity onto the rendered domain,
    \item ternary compatibility:
    \[
    \Pi(x \otimes y \otimes z) = \Pi(x) \otimes_3 \Pi(y) \otimes_3 \Pi(z).
    \]
\end{enumerate}

\subsection*{R8.4 Resonance Visualization Map}

Define the resonance visualization map:
\[
\mathcal{V}_{\alpha}(x,t) = \Pi(\Phi_{\alpha}(x,t)).
\]

This map produces the rendered geometry.

\subsection*{R8.5 Fixed-Point Collapse Under Interpretation}

From Appendix R6, the interpretation operator satisfies:
\[
I_p(\Phi_{\alpha}(x,t)) = c.
\]

Thus the resonance flow collapses to the fixed point under interpretation.

However, the projection $\Pi$ is applied \emph{before} interpretation in
visualization engines, so the rendered geometry is:
\[
\mathcal{V}_{\alpha}(x,t) = \Pi(\Phi_{\alpha}(x,t)),
\]
which is nontrivial.

\subsection*{R8.6 Stability of Resonance Geometry}

Let $A \in \mathcal{O}(X_{42})$ be any admissible operator.

Then:
\[
\Phi_{\alpha}(A(x),t) = A(\Phi_{\alpha}(x,t)).
\]

Applying $\Pi$:
\[
\mathcal{V}_{\alpha}(A(x),t)
= \Pi(A(\Phi_{\alpha}(x,t)))
= \Pi(\Phi_{\alpha}(x,t)).
\]

Thus resonance geometry is invariant under all NDH operators.

\subsection*{R8.7 Oscillatory Behavior}

Define the resonance frequency:
\[
\omega_{\alpha}(x) = \| r_{\alpha}(x) \|.
\]

Since $r_{\alpha}$ is bounded:
\[
0 < \omega_{\alpha}(x) < K,
\]
the resonance flow produces bounded oscillatory motion in $V_3$.

This yields:

\begin{itemize}
    \item pulsing nodes,
    \item coherent oscillation,
    \item stable attractors,
    \item energy spreading,
    \item constellation formation.
\end{itemize}

\subsection*{R8.8 Ternary Collapse and Visualization}

Ternary visualization engines apply a collapse operator:
\[
C_3 : V_3 \to V_3,
\]
satisfying:
\[
C_3(u \otimes_3 v \otimes_3 w) = C_3(u).
\]

Thus:
\[
C_3(\mathcal{V}_{\alpha}(x,t)) = C_3(\Pi(\Phi_{\alpha}(x,t)))
\]
produces stable ternary resonance artifacts.

\subsection*{R8.9 Summary}

The NDH 42D interpretation manifold supports smooth, bounded, ternary-
closed resonance fields. Their flows, when projected into ternary
visualization spaces, produce stable oscillatory structures. These
structures appear as pulsing spheres, shimmering nodes, and coherent
constellations. Interpretation collapses the resonance to the fixed point
$c$, but visualization occurs prior to collapse, yielding the observed
geometry.

This completes Appendix R8.
